\documentclass[conference]{IEEEtran}

\usepackage{booktabs}
\usepackage{graphicx}
\usepackage{amsmath}
\usepackage{amssymb}
\usepackage[T1]{fontenc}
\usepackage{newtxtext}
\usepackage{newtxmath}
\usepackage{microtype}
\usepackage{xcolor}
\usepackage{url}

\setlength{\textfloatsep}{6pt plus 1pt minus 2pt}
\setlength{\floatsep}{6pt plus 1pt minus 2pt}
\setlength{\intextsep}{6pt plus 1pt minus 2pt}

\newcommand{\CorrectedLongTraceTokens}{8{,}192}
\newcommand{\CorrectedHeldoutTraceTokens}{1{,}024}
\newcommand{\CorrectedHeldoutRunCount}{6}
\newcommand{\CorrectedHeldoutSeedBlocks}{3}
\newcommand{\CorrectedHeldoutPairedConditions}{6}
\newcommand{\CorrectedSeedHex}{0xFB72}
\newcommand{\CorrectedMxfpFourLogicalBitsPerValue}{4.25}
\newcommand{\CorrectedSyntheticInitialStateStdDev}{0.05}
\newcommand{\CorrectedSyntheticValueStdDev}{0.25}
\newcommand{\CorrectedRelativeLTwoDenominatorFloor}{1\times 10^{-12}}
\newcommand{\CorrectedNominalAlphaMinimum}{0.950}
\newcommand{\CorrectedNominalAlphaMaximum}{1.000}
\newcommand{\CorrectedNominalBetaMinimum}{0.000}
\newcommand{\CorrectedNominalBetaMaximum}{1.000}
\newcommand{\CorrectedHighRetentionAlphaMinimum}{0.995}
\newcommand{\CorrectedHighRetentionAlphaMaximum}{1.000}
\newcommand{\CorrectedHighRetentionBetaMinimum}{0.850}
\newcommand{\CorrectedHighRetentionBetaMaximum}{1.000}
\newcommand{\CorrectedQdqMxfpFourCosineEightK}{0.729275}
\newcommand{\CorrectedQdqMxfpFourStateRelLTwoEightK}{1.103841}
\newcommand{\CorrectedNativeEncodedCosineEightK}{0.429834}
\newcommand{\CorrectedNativeEncodedStateRelLTwoEightK}{2.580730}
\newcommand{\CorrectedNativeEncodedAlignmentUnderflowsEightK}{2{,}685{,}620{,}693}
\newcommand{\CorrectedBfSixteenCosineEightK}{0.999955}
\newcommand{\CorrectedBfSixteenStateRelLTwoEightK}{0.008988}
\newcommand{\CorrectedDynamicRangeBfSixteenCosineEightK}{0.999875}
\newcommand{\CorrectedDynamicRangeBfSixteenStateRelLTwoEightK}{0.016476}
\newcommand{\CorrectedDynamicRangeQdqMxfpFourCosineEightK}{0.410379}
\newcommand{\CorrectedDynamicRangeQdqMxfpFourStateRelLTwoEightK}{26.772181}
\newcommand{\CorrectedDynamicRangeMxfpEightStateCosineEightK}{0.698070}
\newcommand{\CorrectedDynamicRangeMxfpEightStateRelLTwoEightK}{2.314652}
\newcommand{\CorrectedCancellationBfSixteenCosineEightK}{0.999997}
\newcommand{\CorrectedCancellationBfSixteenStateRelLTwoEightK}{0.008844}
\newcommand{\CorrectedCancellationQdqMxfpFourCosineEightK}{0.981307}
\newcommand{\CorrectedCancellationQdqMxfpFourStateRelLTwoEightK}{1.667123}
\newcommand{\CorrectedCancellationMxfpEightStateCosineEightK}{0.997011}
\newcommand{\CorrectedCancellationMxfpEightStateRelLTwoEightK}{0.255859}
\newcommand{\CorrectedUniformToBfSixteenLutRatio}{1.636}
\newcommand{\CorrectedUniformToBfSixteenStepRatio}{1.437}
\newcommand{\CorrectedHeldoutMinCheckpointCosine}{0.994892}
\newcommand{\CorrectedHeldoutMinAllTokenCosine}{0.994430}
\newcommand{\CorrectedHeldoutMaxFinalStateRelLTwo}{0.097368}
\newcommand{\CorrectedHeldoutMaxAllTokenStateRelLTwo}{0.098653}
\newcommand{\CorrectedTotalHardEvents}{0}
\newcommand{\CorrectedLogicalStateBytes}{537{,}744}
\newcommand{\CorrectedMxfpEightStateBytes}{540{,}672}
\newcommand{\CorrectedBytesBelowMxfpEight}{2{,}928}
\newcommand{\CorrectedPercentBelowMxfpEight}{0.542}
\newcommand{\CorrectedBfSixteenAllLayerStateMiB}{36.00}
\newcommand{\CorrectedUniformMxfpFourAllLayerStateMiB}{9.56}
\newcommand{\CorrectedBfSixteenIdealUram}{1{,}024}
\newcommand{\CorrectedBfSixteenPhysicalFit}{PASS}
\newcommand{\CorrectedBfSixteenRoute}{PASS}
\newcommand{\CorrectedBfSixteenTargetTiming}{FAIL}
\newcommand{\CorrectedBfSixteenTargetWns}{-3.020}
\newcommand{\CorrectedBfSixteenClosingPeriod}{7.125}
\newcommand{\CorrectedBfSixteenClosingFrequency}{140.35}
\newcommand{\CorrectedBfSixteenRoutedLut}{57{,}942}
\newcommand{\CorrectedBfSixteenRoutedFf}{38{,}550}
\newcommand{\CorrectedBfSixteenRoutedBram}{1602.5}
\newcommand{\CorrectedBfSixteenRoutedUram}{928}
\newcommand{\CorrectedBfSixteenRoutedDsp}{15}
\newcommand{\CorrectedBfSixteenVectorlessPower}{5.648}
\newcommand{\CorrectedUniformMxfpFourIdealUram}{256}
\newcommand{\CorrectedUniformMxfpFourIdealBram}{256}
\newcommand{\CorrectedDeviceAvailableUram}{960}
\newcommand{\CorrectedBfSixteenStepInputBytes}{16{,}512}
\newcommand{\CorrectedUniformMxfpFourStepInputBytes}{4{,}480}
\newcommand{\CorrectedCandidateStepInputBytes}{8{,}832}
\newcommand{\CorrectedBfSixteenOutputBytes}{8{,}192}
\newcommand{\CorrectedNativeExpandedOutputBytes}{24{,}576}
\newcommand{\CorrectedBfSixteenLut}{84{,}919}
\newcommand{\CorrectedNativeEncodedMxfpFourLut}{138{,}969}
\newcommand{\CorrectedBfSixteenStepMax}{4{,}225{,}728}
\newcommand{\CorrectedNativeEncodedMxfpFourStepMax}{6{,}070{,}592}
\newcommand{\CorrectedNativeMxfpEightLut}{143{,}084}
\newcommand{\CorrectedNativeMxfpEightStepMax}{6{,}287{,}680}
\newcommand{\CorrectedNativeMxfpEightToBfSixteenLutRatio}{1.685}
\newcommand{\CorrectedNativeMxfpEightToBfSixteenStepRatio}{1.488}
\newcommand{\CorrectedCandidateLut}{435{,}089}
\newcommand{\CorrectedHlsToolVersion}{2025.2}
\newcommand{\CorrectedCandidateNonfoldStepMax}{17{,}954{,}144}
\newcommand{\CorrectedCandidateAmortizedStepMax}{39{,}947{,}365}
\newcommand{\CorrectedCandidateToBfSixteenLutRatio}{5.124}
\newcommand{\CorrectedCandidateToBfSixteenNonfoldRatio}{4.249}
\newcommand{\CorrectedCandidateToBfSixteenAmortizedRatio}{9.453}
\newcommand{\CorrectedCandidatePathNs}{3.108}
\newcommand{\CorrectedConfiguredTimingBudgetNs}{2.920}
\newcommand{\CorrectedConfiguredTimingShortfallNs}{0.188}
\newcommand{\CorrectedCandidateDsp}{78}
\newcommand{\CorrectedPostroutePhysicalFit}{PASS}
\newcommand{\CorrectedPostrouteTwoFiftyStatus}{FAIL}
\newcommand{\CorrectedPostrouteTargetMhz}{250}
\newcommand{\CorrectedPostrouteTwoFiftyWnsNs}{-1.540}
\newcommand{\CorrectedPostrouteTwoHundredStatus}{FAIL}
\newcommand{\CorrectedPostrouteSecondaryMhz}{200}
\newcommand{\CorrectedPostrouteFirstClosingPeriodNs}{5.55}
\newcommand{\CorrectedPostrouteFirstClosingMhz}{180.18}
\newcommand{\CorrectedPostrouteCriticalPathNs}{4.971}
\newcommand{\CorrectedPostrouteCriticalRoutePercent}{92.5}
\newcommand{\CorrectedPostrouteCriticalLogicLevels}{5}
\newcommand{\CorrectedPostrouteLut}{208{,}523}
\newcommand{\CorrectedPostrouteFf}{111{,}027}
\newcommand{\CorrectedPostrouteBram}{353}
\newcommand{\CorrectedPostrouteUram}{624}
\newcommand{\CorrectedPostrouteDsp}{44}
\newcommand{\CorrectedPostrouteUramPercent}{65.0}
\newcommand{\CorrectedPostrouteDrcWarnings}{38}
\newcommand{\CorrectedPostroutePowerTotalW}{6.991}
\newcommand{\CorrectedPostroutePowerPeriodNs}{5.6}
\newcommand{\CorrectedPostroutePowerDynamicW}{3.582}
\newcommand{\CorrectedPostroutePowerStaticW}{3.410}
\newcommand{\CorrectedPostroutePowerConfidence}{Medium}
\newcommand{\CorrectedEZeroControlCosimStatus}{PASS}
\newcommand{\CorrectedEZeroControlCosimCommands}{2}
\newcommand{\CorrectedEZeroControlCosimLatency}{19{,}971}
\newcommand{\CorrectedEZeroRtlSixtyFourStatus}{NOT\_ESTABLISHED}
\newcommand{\CorrectedEZeroRtlTraceTokens}{64}
\newcommand{\CorrectedEZeroRtlTransactions}{1}
\newcommand{\CorrectedEZeroRtlOutputValues}{0}
\newcommand{\CorrectedEZeroHlsCsimSixtyFourStatus}{PASS}
\newcommand{\CorrectedEZeroRtlDirectLoadStatus}{PASS}
\newcommand{\CorrectedEZeroRtlDirectLoadCycles}{4{,}797{,}322}
\newcommand{\CorrectedEZeroRtlCompletedSteps}{0}
\newcommand{\CorrectedMxfpEightPostrouteFit}{PASS}
\newcommand{\CorrectedMxfpEightPostrouteTwoFiftyStatus}{FAIL}
\newcommand{\CorrectedMxfpEightPostrouteTwoFiftyWns}{-1.355}
\newcommand{\CorrectedMxfpEightPostrouteFirstClosingPeriodNs}{6.00}
\newcommand{\CorrectedMxfpEightPostrouteFirstClosingMhz}{166.67}
\newcommand{\CorrectedMxfpEightPostrouteLut}{58{,}582}
\newcommand{\CorrectedMxfpEightPostrouteFf}{39{,}807}
\newcommand{\CorrectedMxfpEightPostrouteBram}{164}
\newcommand{\CorrectedMxfpEightPostrouteUram}{576}
\newcommand{\CorrectedMxfpEightPostrouteDsp}{2}
\newcommand{\CorrectedMxfpEightPostroutePowerTotalW}{4.952}
\newcommand{\CorrectedMxfpEightPostroutePowerPeriodNs}{6.00}
\newcommand{\CorrectedMxfpEightPostroutePowerConfidence}{Medium}
\newcommand{\CorrectedQwenCaptureLayer}{12}
\newcommand{\CorrectedQwenCaptureElements}{147{,}456}
\newcommand{\CorrectedQwenCaptureAbsPnn}{0.6836}
\newcommand{\CorrectedQwenCaptureAbsPnnn}{2.5781}
\newcommand{\CorrectedQwenCaptureAbsMax}{51.25}
\newcommand{\CorrectedQwenCaptureMaxToPnn}{75.0}
\newcommand{\CorrectedQwenRecurrentPromptCount}{4}
\newcommand{\CorrectedQwenRecurrentTotalTokens}{60}
\newcommand{\CorrectedQwenRecurrentMinTokens}{12}
\newcommand{\CorrectedQwenRecurrentMaxTokens}{18}
\newcommand{\CorrectedQwenRecurrentBfSixteenMeanFinalCosine}{0.999991}
\newcommand{\CorrectedQwenRecurrentBfSixteenMeanStateRelLTwo}{0.005894}
\newcommand{\CorrectedQwenRecurrentMxfpFourMeanFinalCosine}{0.982048}
\newcommand{\CorrectedQwenRecurrentMxfpFourMeanStateRelLTwo}{0.324788}
\newcommand{\CorrectedQwenRecurrentMxfpEightStateMeanFinalCosine}{0.989173}
\newcommand{\CorrectedQwenRecurrentMxfpEightStateMeanStateRelLTwo}{0.207192}
\newcommand{\CorrectedQwenRecurrentIntFourMeanFinalCosine}{0.866071}
\newcommand{\CorrectedQwenRecurrentIntFourMeanStateRelLTwo}{0.774523}
\newcommand{\CorrectedQwenRecurrentRsTwoMeanFinalCosine}{0.999635}
\newcommand{\CorrectedQwenRecurrentRsTwoWorstCosine}{0.999408}
\newcommand{\CorrectedQwenRecurrentRsTwoMeanStateRelLTwo}{0.047009}
\newcommand{\CorrectedQwenRecurrentRsTwoWorstStateRelLTwo}{0.050485}
\newcommand{\CorrectedQwenRecurrentRsTwoAccumulatorSaturations}{2}
\newcommand{\CorrectedCandidateFailedIiLoops}{3}
\newcommand{\CorrectedTargetIi}{1}
\newcommand{\CorrectedFailedLoopIi}{2}
\newcommand{\CorrectedCandidateSlrLutPercent}{100}
\newcommand{\CorrectedFoldPeriodTokens}{7}
\newcommand{\CorrectedFoldBurstCyclesMax}{153{,}952{,}544}
\newcommand{\CorrectedTopCommandCyclesMax}{172{,}418{,}717}
\newcommand{\CorrectedCorrectionMachineryLut}{342{,}492}
\newcommand{\CorrectedCorrectionMachineryLutPercent}{78.7}
\newcommand{\CorrectedArithmeticCsimCases}{267}
\newcommand{\CorrectedStatefulCsimTokens}{64}
\newcommand{\CorrectedFixedScaleCosineEightK}{0.847089}
\newcommand{\CorrectedFixedScaleClipsEightK}{67{,}980{,}414}
\newcommand{\CorrectedEveryTokenScaleClipsEightK}{73}
\newcommand{\CorrectedScalePolicyCount}{6}
\newcommand{\CorrectedNumLayers}{36}
\newcommand{\CorrectedNumQkHeads}{16}
\newcommand{\CorrectedNumValueHeads}{32}
\newcommand{\CorrectedKeyDim}{128}
\newcommand{\CorrectedValueDim}{128}
\newcommand{\CorrectedPk}{16}
\newcommand{\CorrectedPv}{8}
\newcommand{\CorrectedBlockSize}{32}
\newcommand{\CorrectedTargetClockNs}{4.0}
\newcommand{\CorrectedLogCapacity}{7}
\newcommand{\CorrectedAccumulatorBits}{32}
\newcommand{\CorrectedAlignmentGuardBits}{5}
\newcommand{\CorrectedEZeroResidualBits}{3}
\newcommand{\CorrectedResidualStackTerms}{2}
\newcommand{\CorrectedOutputCosineThreshold}{0.99}
\newcommand{\CorrectedFinalStateRelLTwoThreshold}{0.10}
\newcommand{\CorrectedStateRelLTwoSensitivityProbe}{0.09}
\newcommand{\CorrectedExtendedStateRelLTwoGateMargin}{0.000943}
\newcommand{\CorrectedCheckpointSet}{64, 256, 1{,}024, 4{,}096, 8{,}192}
\newcommand{\CorrectedETwoMOneSignBits}{1}
\newcommand{\CorrectedETwoMOneExponentBits}{2}
\newcommand{\CorrectedETwoMOneMantissaBits}{1}
\newcommand{\CorrectedETwoMOneMagnitudeCodebook}{0, 0.5, 1, 1.5, 2, 3, 4, 6}
\newcommand{\CorrectedEEightMZeroValidMinimum}{0}
\newcommand{\CorrectedEEightMZeroValidMaximum}{254}
\newcommand{\CorrectedEEightMZeroInvalidCode}{255}
\newcommand{\CorrectedEEightMZeroBias}{127}
\newcommand{\CorrectedEEightMZeroZeroBlockScale}{127}
\newcommand{\CorrectedExtendedQualityStatus}{PASS}
\newcommand{\CorrectedExtendedReplayStatus}{PASS}
\newcommand{\CorrectedExtendedRunCount}{2}
\newcommand{\CorrectedExtendedRecomputedRunCount}{2}
\newcommand{\CorrectedExtendedMinCheckpointCosine}{0.995065}
\newcommand{\CorrectedExtendedMaxFinalStateRelLTwo}{0.097472}
\newcommand{\CorrectedExtendedMinAllTokenCosine}{0.994389}
\newcommand{\CorrectedExtendedMaxAllTokenStateRelLTwo}{0.099057}
\newcommand{\CorrectedExtendedMaxAllTokenStateAbsError}{0.267635}
\newcommand{\CorrectedExtendedTotalElementSaturations}{0}
\newcommand{\CorrectedExtendedTotalAccumulatorSaturations}{0}
\newcommand{\CorrectedExtendedTotalScaleClamps}{0}
\newcommand{\CorrectedExtendedTotalAlignmentUnderflows}{1{,}049{,}355{,}823}
\newcommand{\CorrectedExtendedTotalResidualClips}{23{,}296{,}266}
\newcommand{\CorrectedExtendedTotalFolds}{2{,}340}
\newcommand{\RsTwoLongTokens}{8{,}192}
\newcommand{\RsTwoCheckpointSet}{64, 256, 1{,}024, 4{,}096, and 8{,}192}
\newcommand{\RsTwoBfSixteenCosineEightK}{0.999812}
\newcommand{\RsTwoBfSixteenStateLTwoEightK}{0.020998}
\newcommand{\RsTwoQdqCosineEightK}{0.672818}
\newcommand{\RsTwoQdqStateLTwoEightK}{0.903408}
\newcommand{\RsTwoMxfpEightCosineEightK}{0.950258}
\newcommand{\RsTwoMxfpEightStateLTwoEightK}{0.301757}
\newcommand{\RsTwoIntFourCosineEightK}{0.121427}
\newcommand{\RsTwoIntFourStateLTwoEightK}{2.736632}
\newcommand{\RsTwoCandidateCosineEightK}{0.996725}
\newcommand{\RsTwoCandidateStateLTwoEightK}{0.081526}
\newcommand{\RsTwoCandidateStateMaxAbsEightK}{0.089080}
\newcommand{\RsTwoHeldoutRuns}{6}
\newcommand{\RsTwoHeldoutMinCosine}{0.996099}
\newcommand{\RsTwoHeldoutMaxStateLTwo}{0.082666}
\newcommand{\RsTwoExtendedRuns}{2}
\newcommand{\RsTwoExtendedMinCosine}{0.995974}
\newcommand{\RsTwoExtendedMaxStateLTwo}{0.082737}
\newcommand{\RsTwoExtendedElementSaturations}{0}
\newcommand{\RsTwoExtendedAccumulatorSaturations}{0}
\newcommand{\RsTwoExtendedScaleClamps}{0}
\newcommand{\RsTwoLogicalStateBytes}{576{,}912}
\newcommand{\RsTwoStatePercentOfBfSixteen}{55.02}
\newcommand{\RsTwoStatePercentAboveMxfpEight}{6.70}
\newcommand{\RsTwoAllLayerStateMiB}{19.81}
\newcommand{\RsTwoMxfpEightAllLayerStateMiB}{18.56}
\newcommand{\RsTwoHlsLut}{167{,}082}
\newcommand{\RsTwoHlsDsp}{26}
\newcommand{\RsTwoHlsFmax}{321.96}
\newcommand{\RsTwoHlsNonfoldCycles}{13{,}522{,}144}
\newcommand{\RsTwoHlsFoldCycles}{88{,}629{,}408}
\newcommand{\RsTwoHlsAmortizedCycles}{43{,}065{,}280}
\newcommand{\RsTwoHlsLutRatio}{1.958}
\newcommand{\RsTwoHlsNonfoldRatio}{3.200}
\newcommand{\RsTwoHlsAmortizedRatio}{10.191}
\newcommand{\RsTwoHlsIiStatus}{PASS}
\newcommand{\RsTwoRouteLut}{82{,}038}
\newcommand{\RsTwoRouteFf}{69{,}800}
\newcommand{\RsTwoRouteBram}{341}
\newcommand{\RsTwoRouteUram}{598}
\newcommand{\RsTwoRouteDsp}{12}
\newcommand{\RsTwoRouteTargetWns}{-1.736}
\newcommand{\RsTwoRouteTargetMhz}{250}
\newcommand{\RsTwoRouteClosingPeriodNs}{6.0}
\newcommand{\RsTwoRouteClosingMhz}{166.67}
\newcommand{\RsTwoRoutePowerW}{5.242}
\newcommand{\RsTwoRouteDynamicW}{1.874}
\newcommand{\RsTwoRouteStaticW}{3.368}
\newcommand{\RsTwoRtlStatus}{PASS}
\newcommand{\RsTwoRtlTokens}{64}
\newcommand{\RsTwoRtlTransactions}{66}
\newcommand{\RsTwoRtlOutputs}{262{,}144}
\newcommand{\RsTwoRtlStepMean}{44{,}824{,}405}


\title{Can Native MXFP4 Replace BF16 in FPGA Gated DeltaNet Decode?\\
Cost and Long-Horizon Recurrent-State Drift}

\author{\IEEEauthorblockN{Anonymous Authors}}

\begin{document}
\maketitle

\begin{abstract}
This paper asks whether native Open Compute Project (OCP) MXFP4 arithmetic can
replace bfloat16 (BF16)-style arithmetic in a persistent-state Gated DeltaNet
recurrence-core kernel on an AMD Alveo U55C field-programmable gate array
(FPGA). Persistent state and the five-phase token schedule are inherited from
prior work and held fixed. We implement E2M1 products in integer lookup-table (LUT) logic and
E8M0-aware accumulation, then separate floating quantize/dequantize (Q/DQ)
diagnostics from the native encoded path. At token \CorrectedLongTraceTokens,
BF16 gives \CorrectedBfSixteenCosineEightK{} output cosine and
\CorrectedBfSixteenStateRelLTwoEightK{} recurrent-state relative $L_2$ error
against single-precision floating point (FP32). MXFP4 Q/DQ gives
\CorrectedQdqMxfpFourCosineEightK{} and
\CorrectedQdqMxfpFourStateRelLTwoEightK{}, whereas the native encoded path gives
\CorrectedNativeEncodedCosineEightK{} and
\CorrectedNativeEncodedStateRelLTwoEightK{}. The encoded HLS kernel uses
\CorrectedUniformToBfSixteenLutRatio{} times the LUTs and
\CorrectedUniformToBfSixteenStepRatio{} times the maximum estimated STEP cycles
of the fixed-geometry BF16 kernel. A recurrence-aware residual/write-log
candidate passes \CorrectedHeldoutSeedBlocks{} test seed blocks crossed with
two initial-state conditions and two \CorrectedLongTraceTokens{}-token
development conditions, but uses \CorrectedCandidateToBfSixteenLutRatio{} times
the BF16 LUTs and misses \CorrectedCandidateFailedIiLoops{} explicit II=1
constraints. Thus the present native MXFP4 implementations do not support
replacing BF16. The corrected candidate's out-of-context post-route image fits
using \CorrectedPostrouteLut{} CLB LUTs and \CorrectedPostrouteUram{} of
\CorrectedDeviceAvailableUram{} URAMs, but has
\CorrectedPostrouteTwoFiftyWnsNs{} ns worst negative slack at
\CorrectedPostrouteTargetMhz{} MHz; the first passing point in the tested
fixed-route sweep is \CorrectedPostrouteFirstClosingMhz{} MHz. The long-horizon
quality result is layer-level and synthetic. Four short model-derived Qwen
layer-12 recurrence traces preserve the same arithmetic ordering but contain
only \CorrectedQwenRecurrentMinTokens{}--\CorrectedQwenRecurrentMaxTokens{}
valid tokens each. Power is a vectorless estimate rather than board energy.
\end{abstract}

\begin{IEEEkeywords}
Gated DeltaNet, MXFP4, recurrent state, FPGA, numerical drift, Vitis HLS
\end{IEEEkeywords}

\section{Introduction}
Gated DeltaNet (GDN) replaces a growing attention history with a fixed-size
recurrent matrix that is decayed and updated once per token
\cite{gateddeltanet}. This makes the numerical state persistent: an arithmetic
error written at one token can affect every later token. Lowering state
precision therefore changes more than storage cost. It changes the dynamical
system implemented by the decoder.

Gupta et al. map GDN decode to a field-programmable gate array (FPGA) using persistent on-chip state and a
five-phase BF16 pipeline \cite{gupta2026}. We use that organization as the
baseline dataflow. Neither persistent state nor the five-phase schedule is a
contribution of this work. We study only the arithmetic replacement inside that
recurrence-core boundary.

The Open Compute Project (OCP) MXFP4 format stores \CorrectedBlockSize{} E2M1
elements with one E8M0 scale
\cite{ocpmx}. Its \CorrectedMxfpFourLogicalBitsPerValue{} logical bits per value
are attractive for a recurrent
state that is read and rewritten every token. The shared exponent also creates
two risks. Values below the block range disappear during alignment, and a
scale change requantizes the complete block. Those errors feed back through
the next state update.

We evaluate the replacement as a controlled experiment. The recurrence,
logical K-by-V state, layer dimensions, U55C target,
\CorrectedTargetClockNs{} ns constraint, $P_K=\CorrectedPk$,
$P_V=\CorrectedPv$, and block size \CorrectedBlockSize{} are fixed. The study
contributes:

\begin{enumerate}
  \item an encoded-integer MXFP4 recurrence oracle and native high-level
        synthesis (HLS) arithmetic with explicit E2M1 decoding, E8M0 alignment,
        round-to-nearest-even (RNE), Q1.15 controls, and event counters;
  \item long synthetic traces at \CorrectedCheckpointSet{} tokens for
        FP32, BF16, MXFP4 floating Q/DQ, native encoded MXFP4, flat INT4, and
        a floating-Q/DQ MXFP8-state fallback, plus four short model-derived
        recurrence traces and
        \CorrectedScalePolicyCount{} scale
        refresh policies;
  \item a fixed-geometry BF16, encoded-MXFP4, and native-MXFP8 HLS comparison and a
        separately classified recurrence-aware correction, including
        out-of-context post-route timing, utilization, design-rule, and
        vectorless-power evidence for the correction and MXFP8 baseline; and
  \item an evidence-backed negative answer for the current implementations,
        including the numerical failure of native encoded MXFP4 and the hardware cost
        of the correction that meets the frozen long-horizon synthetic gate.
\end{enumerate}

The evaluation covers one recurrence core, synthetic recurrent traces, four
short model-derived layer-12 recurrence traces, matched HLS estimates, and
out-of-context implementation attempts. It does not measure Qwen3-Next
perplexity, downstream accuracy, or board energy. Within that boundary, can native MXFP4 replace BF16-style
arithmetic while reducing FPGA cost and limiting empirical recurrent-state
drift over long synthetic decode traces?

\section{Background and Scope}
\subsection{Microscaling Representation}
An E2M1 element contains \CorrectedETwoMOneSignBits{} sign bit,
\CorrectedETwoMOneExponentBits{} exponent bits, and
\CorrectedETwoMOneMantissaBits{} mantissa bit. Its positive finite magnitude
codebook is $\{\CorrectedETwoMOneMagnitudeCodebook\}$. Both signed-zero codes
decode to zero, and the encoder emits positive zero. An E8M0 byte encodes the
shared power-of-two scale $2^{e-\CorrectedEEightMZeroBias{}}$ for codes
\CorrectedEEightMZeroValidMinimum{} through
\CorrectedEEightMZeroValidMaximum{}. Code
\CorrectedEEightMZeroInvalidCode{} is invalid; an all-zero block uses scale code
\CorrectedEEightMZeroZeroBlockScale{}. NaN and infinity inputs are rejected.
These rules follow the OCP representation \cite{ocpmx}; scale selection and RNE
at each state-write boundary are implementation choices.

\subsection{Recurrence-Core Boundary}
The head counts and K-by-V dimensions follow the pinned official
Qwen3-Next-80B-A3B-Instruct configuration and Transformers implementation
\cite{qwen3nextmodel,transformersqwen}.
For value head $h$, query/key head $p(h)=\lfloor h/2\rfloor$ serves the
corresponding pair of value heads. With logical state
$S_h\in\mathbb{R}^{K\times V}$ and
$K=\CorrectedKeyDim$, $V=\CorrectedValueDim$, the
token recurrence is
\begin{align}
S_h^{d} &= \alpha_h S_h, \\
u_h &= \beta_h\left(v_h-k_{p(h)}^\mathsf{T}S_h^{d}\right), \\
S_h' &= S_h^{d}+k_{p(h)}u_h^\mathsf{T}, \\
o_h &= q_{p(h)}^\mathsf{T}S_h'.
\end{align}
The query is normalized and scaled by $1/\sqrt{\CorrectedKeyDim}$ before entering the
kernel; the key is normalized. Gated RMSNorm, the $z$ branch, projections,
depthwise convolution, and output projection are outside the measured
boundary. We therefore call the implementation a GDN recurrence-core kernel,
not a complete GDN layer.
Holding the unchanged convolution, projections, and output normalization outside
the boundary isolates arithmetic-induced drift and cost in the persistent
recurrence rather than mixing them with unmodified layer components.

The operation order is decay, prediction from the decayed state, delta
formation, rank-one write, and output from the updated state. The software
oracle was checked against pinned Transformers and FLA recurrent, chunk, and
cache paths before quantized experiments were admitted
\cite{transformersqwen,fla021}.

\subsection{Prior Work and Component Ownership}
DeltaNet establishes the rank-one erase/write recurrence and parallel training
formulations on which GDN builds \cite{deltanet}; the gated recurrence itself is
inherited from Yang et al. \cite{gateddeltanet}. The persistent-state FPGA
organization, paired-head handling, and five phases are inherited from Gupta
et al. \cite{gupta2026}. MXFP4-B32 is standardized \cite{ocpmx}. QLlama maps
microscaling quantization to FPGA inference, while Jack of All Scales studies
MXFP dot-product mappings and digital signal-processing (DSP) block support on
Agilex-5 \cite{qllama,jackscales}. MXFormer implements block-exponent-aligned
MXFP4 compute in charge-trap memory, MicroMix mixes MXFP4/6/8 channels, and
MXAttention and an MXFP post-training-quantization benchmark identify scaling,
clipping, and underflow as material quality factors
\cite{mxformer,micromix,mxattention,mxfpbenchmark}. These are precedents for
alignment, mixed precision, and scale-policy studies. Zhang et al. quantize the chunk-wise GDN
matrix-inversion path rather than token-serial recurrent state
\cite{quantgdn}. KIVI retains recent full-precision entries beside a quantized
cache, and SSDi8 uses a persistent INT8 state-space path with error correction
\cite{kivi,ssdi8}. These works establish prior art for the component ideas. Our
scope is the matched arithmetic-replacement experiment and its verification,
not the invention of those constituents.
Table~\ref{tab:ownership} makes the resulting ownership boundaries explicit.

\begin{table}[t]
\centering
\caption{Component ownership relative to the corrected implementation.}
\label{tab:ownership}
\footnotesize
\begin{tabular}{@{}p{0.36\linewidth}p{0.18\linewidth}p{0.33\linewidth}@{}}
\toprule
Component & Classification & Basis \\
\midrule
Persistent on-chip state & Inherited & Gupta et al. \cite{gupta2026} \\
Five-phase token schedule & Inherited & Gupta et al. \cite{gupta2026} \\
Qwen recurrence and core boundary & Corrected & Pinned software parity \\
OCP MXFP4-B32 state & Modified & Standard format in a recurrent path \\
E2M1/E8M0 FPGA arithmetic & Modified & OCP format plus QLlama and Jack of All
Scales \cite{ocpmx,qllama,jackscales} \\
Encoded arithmetic oracle & Newly proposed & Artifact-specific independent
integer implementation \\
E2M1/E2M0 residual plus fixed write log & Modified & Related
residual-buffer methods exist \cite{kivi,ssdi8}; no first-use claim \\
Long-horizon drift harness & Newly proposed & Same recurrence and synthetic
input prefix \\
\addlinespace
\multicolumn{3}{p{0.92\linewidth}}{\emph{Note:} Newly proposed identifies
ownership within this artifact, not a claim of literature priority.} \\
\bottomrule
\end{tabular}
\end{table}

\section{Numerical Method}
\subsection{Uniform MXFP4 Path}
The matched BF16 comparator rounds state and operands to BF16 at their declared
boundaries and performs reductions in FP32. The MXFP8-state fallback changes
only the recurrent state to eight-bit E4M3 elements under E8M0 B32 scales while
retaining MXFP4 token operands; flat INT4 uses signed symmetric four-bit quantization
with a tensor-wide scale.

We report a floating Q/DQ MXFP4 diagnostic and a distinct native encoded path.
The Q/DQ diagnostic quantizes at named boundaries but performs the intervening
recurrence in floating point. It does not represent the HLS accumulator or its
event behavior. The encoded path below is the software counterpart of the
native HLS arithmetic.

The uniform path stores E2M1 elements in B32 blocks with an E8M0 scale per
block. For nonzero block maximum $m$ and largest finite E2M1 magnitude
$x_{\max}$, the selected scale power is
$s=\lceil\log_2(m/x_{\max})\rceil$; its biased E8M0 code is clamped to the
valid range. Elements are divided by $2^s$ and rounded to E2M1 with RNE. Zero
has a canonical encoding. If E2M1 magnitude code $r$ denotes $r/2$, a product
has integer magnitude $r_a r_b$ and power $s_a+s_b-2$.
One flat 128-term reduction aligns every product directly to the greatest nonzero product exponent (power),
applies RNE to the full exponent-difference shift, and accumulates in ascending
key-index order with signed \CorrectedAccumulatorBits{}-bit saturation after
every addition. A nonzero term rounded to zero is counted as an alignment
underflow. This uniform baseline has no guard-bit offset; the five guard bits
described below belong only to the corrected candidate.

Each synthetic STEP record supplies normalized $q$ and $k$, value $v$, decay
$\alpha$, delta gate $\beta$, and the selected initial state $S$. The MXFP4
paths encode $q$, $k$, $v$, and every state write; $\alpha$ and $\beta$ use
unsigned Q1.15 fixed point. There are no projection weights in a test vector
because all projections are outside the recurrence-core boundary. The MXFP8
fallback changes only state element codes to E4M3 under B32 E8M0 scales; flat
INT4 changes the same token operands and state under its declared tensor-wide
scale. Requantization, element and accumulator saturation, scale clamps,
alignment underflows, and scale changes are explicit state-write events. The
complete contract is versioned in the artifact rather than inferred from C++
types.

The element product uses sign exclusive-or (XOR), table-based E2M1 magnitude
multiplication, and exponent addition. No floating-point type appears in the
active native arithmetic sources. The resulting E2M1 element-product logic is
integer based;
wider coefficient operations may still map to DSPs and are reported as such.

\subsection{Recurrence-Aware Correction}
Native encoded MXFP4 fails the frozen finite-trace thresholds, so we also
evaluate one locally frozen correction. The base state has an E2M1 primary term and a dense
signed \CorrectedEZeroResidualBits{}-bit E2M0-style residual, each with B32
scales. \CorrectedResidualStackTerms{}-term residual-stacked E2M1 vectors
represent token inputs and live rank-one writes.
Up to \CorrectedLogCapacity{} writes remain in a log with Q1.15 decay
coefficients. At capacity, all \CorrectedLogCapacity{} entries are folded
tilewise into the base and the log is cleared.
This candidate aligns shared exponents with \CorrectedAlignmentGuardBits{}
guard bits before signed \CorrectedAccumulatorBits{}-bit accumulation.

This representation is a mitigation, not the controlled uniform replacement.
It changes the physical state representation and adds a periodic fold. Its
logical payload is \CorrectedLogicalStateBytes{} bytes per layer, only
\CorrectedBytesBelowMxfpEight{} bytes below a uniform
MXFP8-E4M3-B32 state. Exact arithmetic equivalence was tested separately before
quantized evaluation.

\section{FPGA Implementation}
Both controlled baselines use equivalent persistent-state command semantics and
the same five-phase token organization. Commands select one of
\CorrectedNumLayers{} runtime layer slots and support RESET, LOAD, STEP, and
READBACK. Their data-plane encodings are not identical: BF16 transports BF16
tensors, whereas MXFP4 transports encoded elements and scales and returns
expanded integer output fields. Packing and token quantization are outside both
HLS kernels. Before interface packing, a STEP consumes
\CorrectedBfSixteenStepInputBytes{} logical input bytes in the BF16 kernel and
\CorrectedUniformMxfpFourStepInputBytes{} in the uniform MXFP4 kernel; the
corrected residual-stack candidate consumes
\CorrectedCandidateStepInputBytes{}. BF16 returns
\CorrectedBfSixteenOutputBytes{} logical output bytes, while the native kernels'
expanded INT32-mantissa/INT16-exponent fields total
\CorrectedNativeExpandedOutputBytes{}. These are type-layout counts, not
measured AXI traffic; RESET has no tensor payload, and resident-state LOAD and
READBACK use their separately reported logical state payloads. The HLS target is
\texttt{xcu55c-fsvh2892-2L-e} under Vitis HLS
\CorrectedHlsToolVersion{} at \CorrectedTargetClockNs{} ns with
$P_K=\CorrectedPk$, $P_V=\CorrectedPv$, and B32.
We retain the U55C target used by the inherited BF16 dataflow so the arithmetic
comparison does not also change devices. We fix this target and parallelism
point for the replacement rather than tune each format independently.

The corrected candidate adds packed primary and residual banks,
\CorrectedLogCapacity{} packed key/update log entries, per-head gamma and lambda
coefficients, and event counters. A non-mutating snapshot exposes all
persistent fields for exact verification. Folding is atomic at the command
boundary.

Figure~\ref{fig:datapath} maps the correction onto the inherited five-phase
schedule. The base state remains resident while each STEP decays its
coefficient, evaluates the base and live writes, appends the new rank-one
update, and reads the output from the updated logical state. When the
\CorrectedLogCapacity{}-entry log fills, a tilewise fold runs after the output
has observed the newest write.

\begin{figure*}[t]
\centering
\includegraphics[width=\textwidth]{paper/figures/corrected/corrected_candidate_datapath.pdf}
\caption{Corrected-candidate datapath on the inherited five-phase token
schedule. Gray denotes inherited control and scheduling, orange denotes native
MX arithmetic, teal and blue denote the corrected persistent representation,
and the dashed path is the periodic atomic fold. Dimensions and rates are
generated from the numerical contract; the drawing is logical architecture,
not physical-placement evidence.}
\label{fig:datapath}
\end{figure*}

Logically, B32 blocks run along the V-column axis within each K row. The code
partitions the K and V access lanes by $P_K$ and $P_V$ and targets an initiation
interval (II) of \CorrectedTargetIi{} on inner loops. Source directives request
block RAM and UltraRAM bindings, but the inferred memory counts contradict
bit-capacity lower bounds for the declared deep arrays and are excluded from
paper conclusions.
After removing invalid pipeline directives from long validation and
requantization loops, HLS estimates a \CorrectedCandidatePathNs{} ns path. The
\CorrectedTargetClockNs{} ns target and tool-reported uncertainty imply a
\CorrectedConfiguredTimingBudgetNs{} ns effective budget, so the estimate
misses the configured margin by \CorrectedConfiguredTimingShortfallNs{} ns.
\CorrectedCandidateFailedIiLoops{} remaining targeted loops achieve
II=\CorrectedFailedLoopIi{} rather than II=\CorrectedTargetIi{}. HLS reports
\CorrectedCandidateDsp{} DSPs for coefficient and
wider arithmetic; the E2M1 element multiplier itself remains integer logic.
These source bindings and estimates do not establish a physically allocated
\CorrectedNumLayers{}-layer bank by themselves;
Section~\ref{sec:postroute} separately reports the corrected candidate's
out-of-context implementation.

\section{Evaluation}
\subsection{Questions and Controls}
The experiments answer three questions:

\begin{itemize}
  \item Does native encoded MXFP4 limit empirical output and recurrent-state
        drift over
        \CorrectedLongTraceTokens{} synthetic tokens relative to FP32 and BF16?
  \item Do floating-Q/DQ scale refresh policies or an MXFP8-E4M3 state
        fallback meet the same finite-trace criteria?
  \item If a recurrence-aware correction meets the finite-trace thresholds,
        does its HLS estimate improve on the fixed-geometry BF16 kernel?
  \item Does that corrected candidate physically fit and close the declared
        clock target after out-of-context routing?
\end{itemize}

The controlled nominal baseline variants use \CorrectedNumQkHeads{} query/key heads,
\CorrectedNumValueHeads{} value heads, $K=\CorrectedKeyDim$,
$V=\CorrectedValueDim$, the
same K-by-V state, and the same generated input prefix. The development seed is
\texttt{\CorrectedSeedHex}. Required checkpoints are \CorrectedCheckpointSet.
The nominal trace draws query and key vectors from independent standard normal
distributions before normalization. Values and random initial states are
zero-mean normal with standard deviations \CorrectedSyntheticValueStdDev{} and
\CorrectedSyntheticInitialStateStdDev{}, respectively. Nominal decay and delta
gates follow $\mathcal{U}[\CorrectedNominalAlphaMinimum,
\CorrectedNominalAlphaMaximum]$ and
$\mathcal{U}[\CorrectedNominalBetaMinimum,\CorrectedNominalBetaMaximum]$.
The high-retention correction traces keep the vector distributions but use
$\alpha\sim\mathcal{U}[\CorrectedHighRetentionAlphaMinimum,
\CorrectedHighRetentionAlphaMaximum]$ and
$\beta\sim\mathcal{U}[\CorrectedHighRetentionBetaMinimum,
\CorrectedHighRetentionBetaMaximum]$, with both random and zero initial states.
Changing the trace-family tag changes the deterministic RNG stream, even at the
same seed. The correction results are therefore not paired numerical ablations
of the nominal baseline trace.
The frozen engineering gate requires output cosine at least
\CorrectedOutputCosineThreshold{} at every checkpoint, final state relative
$L_2$ at most \CorrectedFinalStateRelLTwoThreshold{}, and zero element
saturation, accumulator saturation, and scale-clamp events. These are the three
hard-event counters; alignment underflows and deliberate E2M0 residual clips
are reported but excluded from the zero-event condition. We fixed these
complementary diagnostics before test-set access: cosine constrains observable
output direction, while state relative $L_2$ constrains hidden-state magnitude.
They were not calibrated to downstream model quality.

The correction uses \CorrectedHeldoutSeedBlocks{} test seed blocks crossed with
random and zero initial states, producing \CorrectedHeldoutPairedConditions{}
paired \CorrectedHeldoutTraceTokens{}-token conditions. Every artifact was
deterministically recomputed by the same implementation and compared field by
field with its recorded comma-separated values (CSV) artifact. This checks
integrity and determinism, not implementation independence. The
\CorrectedExtendedRunCount{} \CorrectedLongTraceTokens{}-token
development runs use the same
recorded seed with random and zero initial states. Each longer run was fully
recomputed from source and compared field by field with its recorded CSV
artifact. The encoded HLS candidate was
also checked for \CorrectedArithmeticCsimCases{} frozen arithmetic cases and for
\CorrectedStatefulCsimTokens{} stateful C-simulation tokens through repeated
folds and final snapshot readback.

The corrected candidate is synthesized and routed out of context for the
declared U55C part. A source-level clock-buffer insertion is hash-checked
against the HLS-generated top before implementation. Timing is re-reported on
the same routed checkpoint at several clock periods; the first passing tested
point is not interpreted as a binary-searched maximum frequency. Post-route
power uses Vivado vectorless activity propagation at a timing-closed
\CorrectedPostroutePowerPeriodNs{} ns period. It is a tool estimate, not board
telemetry. Vitis HLS/XSIM also runs
\CorrectedEZeroControlCosimCommands{} exact early-return commands through the
generated RTL. Those commands exercise control/status behavior but no recurrent
state transition. A direct AXI harness over the same HLS-generated Verilog
additionally completes an exact candidate LOAD in
\CorrectedEZeroRtlDirectLoadCycles{} cycles. Two isolated official 64-token
XSIM attempts exhaust host memory before transaction one, so neither result is
reported as candidate recurrent RTL parity.

Separately, a pinned real Qwen3-Next layer-input capture at layer
\CorrectedQwenCaptureLayer{} is combined with the matching checkpoint
projection weights to reconstruct query, key, value, decay, and delta-gate
inputs at the post-convolution, pre-query/key-normalization recurrence
boundary. The \CorrectedQwenRecurrentPromptCount{} prompt traces contain
\CorrectedQwenRecurrentTotalTokens{} valid tokens in total, with
\CorrectedQwenRecurrentMinTokens{}--\CorrectedQwenRecurrentMaxTokens{} tokens
per prompt. They are evaluated as floating Q/DQ diagnostics and are not mixed
into the synthetic long-horizon curves. This is a short real-input recurrence
check, not closed-loop quantized model inference.

\subsection{Metrics and Evidence Levels}
Output cosine similarity flattens all value heads and V coordinates at each
token. State relative $L_2$ flattens all heads, K rows, and V columns and is
$\|S_q-S_{32}\|_F/\max(\|S_{32}\|_F,
\CorrectedRelativeLTwoDenominatorFloor)$; maximum absolute state error and
arithmetic events are also recorded. Exact equality returns cosine one and zero
error; the artifact specifies the near-zero-norm cosine convention. No headwise
minimum, percentile, or acceptance gate is computed. Software traces, HLS C simulation, HLS
synthesis,
register-transfer-level (RTL) simulation, post-route implementation, and board
measurement are separate evidence levels. A PASS at one level does not fill a
missing later level.

\section{Results}
\subsection{Uniform Arithmetic Replacement}
Figure~\ref{fig:long} and Table~\ref{tab:long} show the controlled synthetic
trace. BF16 remains close to FP32 through the final checkpoint. Both MXFP4 Q/DQ
and native encoded MXFP4 fail the frozen criteria, and the encoded path drifts
more than its floating diagnostic. Flat INT4 also fails. The floating-Q/DQ
MXFP8-E4M3 state fallback
improves on four-bit state but still misses both final criteria. Thus the Q/DQ
curve cannot be paired with the encoded HLS cost as if they were one
implementation.

Two supplemental full-geometry, \CorrectedLongTraceTokens{}-token Q/DQ
traces test deterministic exponent cycling and alternating-sign cancellation.
Each pair below reports (output cosine, state relative $L_2$) at token 8192.
Dynamic range gives BF16
$($\CorrectedDynamicRangeBfSixteenCosineEightK{},
\CorrectedDynamicRangeBfSixteenStateRelLTwoEightK{}$)$, MXFP4
$($\CorrectedDynamicRangeQdqMxfpFourCosineEightK{},
\CorrectedDynamicRangeQdqMxfpFourStateRelLTwoEightK{}$)$, and MXFP8 state
$($\CorrectedDynamicRangeMxfpEightStateCosineEightK{},
\CorrectedDynamicRangeMxfpEightStateRelLTwoEightK{}$)$.
Cancellation gives BF16
$($\CorrectedCancellationBfSixteenCosineEightK{},
\CorrectedCancellationBfSixteenStateRelLTwoEightK{}$)$, MXFP4
$($\CorrectedCancellationQdqMxfpFourCosineEightK{},
\CorrectedCancellationQdqMxfpFourStateRelLTwoEightK{}$)$, and MXFP8 state
$($\CorrectedCancellationMxfpEightStateCosineEightK{},
\CorrectedCancellationMxfpEightStateRelLTwoEightK{}$)$.
BF16 stays within both full-trace thresholds; every low-precision variant
fails at least one, although MXFP8 state is the strongest low-precision
comparator. These traces are synthetic range/cancellation probes, not
model-derived outlier evidence.

\paragraph{Real-capture checks.}
The pinned Qwen layer-input capture contains
\CorrectedQwenCaptureElements{} values. Its absolute-value 99th and 99.9th
percentiles are \CorrectedQwenCaptureAbsPnn{} and
\CorrectedQwenCaptureAbsPnnn{}, while the maximum is
\CorrectedQwenCaptureAbsMax{}, or \CorrectedQwenCaptureMaxToPnn{} times the
99th percentile. This confirms that a real layer boundary can contain a much
wider tail than the nominal Gaussian trace. The reconstructed recurrence traces
in Table~\ref{tab:qwen-recurrence} give BF16 mean final output cosine/state
relative $L_2$ of
\CorrectedQwenRecurrentBfSixteenMeanFinalCosine{}/
\CorrectedQwenRecurrentBfSixteenMeanStateRelLTwo{}, MXFP4 gives
\CorrectedQwenRecurrentMxfpFourMeanFinalCosine{}/
\CorrectedQwenRecurrentMxfpFourMeanStateRelLTwo{}, MXFP8 state gives
\CorrectedQwenRecurrentMxfpEightStateMeanFinalCosine{}/
\CorrectedQwenRecurrentMxfpEightStateMeanStateRelLTwo{}, and INT4 gives
\CorrectedQwenRecurrentIntFourMeanFinalCosine{}/
\CorrectedQwenRecurrentIntFourMeanStateRelLTwo{}. The ordering agrees with the
synthetic study, but these short traces cannot establish long-horizon stability,
perplexity, or downstream accuracy.

\begin{table*}[t]
\centering
\caption{Floating-Q/DQ recurrence diagnostics on four model-derived Qwen
layer-12 prompt traces. Final values are averaged per prompt; ``worst'' and
maximum absolute error range over all prompts and valid token positions.}
\label{tab:qwen-recurrence}
\scriptsize
\resizebox{\textwidth}{!}{\begin{tabular}{lrrrrr@{\hspace{1em}}r}
\toprule
Arithmetic & Mean final cos. & Worst cos. & Mean state $L_2$ & Worst state $L_2$ & Max abs. & Nonfinite \\
\midrule
BF16 state/operands & 0.99999 & 0.99998 & 0.00589 & 0.00657 & 0.01917 & 0 \\
MXFP4 state/operands & 0.98205 & 0.96768 & 0.32479 & 0.33888 & 0.61595 & 0 \\
MXFP4 operands + MXFP8 state & 0.98917 & 0.96967 & 0.20719 & 0.21994 & 0.46683 & 0 \\
Flat INT4 & 0.86607 & 0.75192 & 0.77452 & 0.82253 & 0.40506 & 0 \\
Native MXFP4 RS2/R3 & 0.99963 & 0.99941 & 0.04701 & 0.05049 & 0.07849 & 0 \\
\bottomrule
\end{tabular}
}
\end{table*}

\begin{figure*}[t]
\centering
\includegraphics[width=\textwidth]{paper/figures/corrected/long_sequence_stability_panel.pdf}
\caption{Layer-level synthetic output cosine (left) and recurrent-state relative
$L_2$ error (right) versus token index. The horizontal threshold guides are
engineering criteria, not model-quality guarantees. Each line is one
deterministic nominal development trace. Curves distinguish floating Q/DQ from
native encoded MXFP4; one seed cannot estimate trace-population uncertainty.}
\label{fig:long}
\end{figure*}

\begin{table*}[t]
\centering
\caption{Controlled synthetic drift at every required checkpoint. Each
cell gives output cosine and state relative $L_2$ against FP32.}
\label{tab:long}
\scriptsize
\resizebox{\textwidth}{!}{\begin{tabular}{lrrrrrrrrrr}
\toprule
& \multicolumn{2}{c}{64} & \multicolumn{2}{c}{256} & \multicolumn{2}{c}{1{,}024} & \multicolumn{2}{c}{4{,}096} & \multicolumn{2}{c}{8{,}192} \\
Variant & Cos. & Rel. $L_2$ & Cos. & Rel. $L_2$ & Cos. & Rel. $L_2$ & Cos. & Rel. $L_2$ & Cos. & Rel. $L_2$ \\
\midrule
BF16 state/operands & 0.99996 & 0.00853 & 0.99996 & 0.00905 & 0.99996 & 0.00873 & 0.99996 & 0.00898 & 0.99996 & 0.00899 \\
MXFP4 floating Q/DQ & 0.86488 & 0.69728 & 0.76740 & 1.08234 & 0.72694 & 1.13178 & 0.73492 & 1.13826 & 0.72927 & 1.10384 \\
Native encoded MXFP4 & 0.79393 & 0.98196 & 0.60220 & 1.81363 & 0.45672 & 2.34972 & 0.40960 & 2.62835 & 0.42983 & 2.58073 \\
MXFP4 Q/DQ + MXFP8-E4M3 state & 0.97213 & 0.22072 & 0.97075 & 0.23507 & 0.96922 & 0.22908 & 0.97059 & 0.22996 & 0.96850 & 0.23213 \\
Flat INT4 & 0.63811 & 1.50722 & 0.55497 & 1.80492 & 0.50944 & 1.87392 & 0.55845 & 1.52630 & 0.52540 & 1.86201 \\
\bottomrule
\end{tabular}
}
\end{table*}

No tested floating-Q/DQ scale policy rescues uniform MXFP4. Fixed scales retain
more output direction but produce extensive state-element clipping; frequent
refresh reduces clipping without satisfying the drift gate. Periodic and
threshold-triggered refresh also fail.
Table~\ref{tab:scales} reports the complete token-\CorrectedLongTraceTokens{}
ablation rather than
selecting one favorable
metric.

\begin{table*}[t]
\centering
\caption{Uniform-MXFP4 block-scale policies at token \CorrectedLongTraceTokens.
Clips and scale
changes are cumulative floating quantize/dequantize diagnostics.}
\label{tab:scales}
\footnotesize
\begin{tabular}{lrrrr}
\toprule
Scale policy & Output cosine & State rel. $L_2$ & Clips & Scale changes \\
\midrule
Fixed & 0.847089 & 0.676620 & 67{,}980{,}414 & 0 \\
Every token & 0.729275 & 1.103841 & 73 & 26{,}504{,}191 \\
Periodic, $N=4$ & 0.773418 & 0.985345 & 16{,}091{,}670 & 8{,}345{,}806 \\
Periodic, $N=8$ & 0.785090 & 0.951738 & 20{,}026{,}426 & 4{,}448{,}732 \\
Periodic, $N=16$ & 0.785347 & 0.939046 & 22{,}646{,}994 & 2{,}354{,}947 \\
Threshold triggered & 0.558410 & 1.724816 & 0 & 35{,}057 \\
\bottomrule
\end{tabular}

\end{table*}

\subsection{Fixed-Geometry HLS Cost}
Table~\ref{tab:hls} holds the recurrence, dataflow, target, parallelism, and
block size fixed. The native encoded MXFP4 point is larger and has a longer
estimated STEP schedule than BF16, so the HLS cost hypothesis fails. Native
MXFP8 is also synthesized at the same boundary as a wider low-precision
baseline. The
different data-plane encodings and excluded packing mean this is not an
end-to-end interface or energy ratio. All inferred memory counts are omitted
because they conflict with bit-capacity lower bounds for the declared arrays.

\begin{table}[t]
\centering
\caption{Fixed-geometry U55C HLS logic and schedule estimates. Memory counts,
packing, and transport are excluded.}
\label{tab:hls}
\scriptsize
\resizebox{\columnwidth}{!}{\begin{tabular}{lrrrrr}
\toprule
Arithmetic & Path (ns) & STEP max & LUT & FF & DSP \\
\midrule
BF16 & 2.920 & 4{,}225{,}728 & 84{,}919 & 39{,}087 & 0 \\
Native encoded MXFP4 & 2.920 & 6{,}070{,}592 & 138{,}969 & 43{,}962 & 3 \\
Native MXFP8 & 2.920 & 6{,}287{,}680 & 143{,}084 & 44{,}715 & 3 \\
\bottomrule
\end{tabular}
}
\end{table}

Figure~\ref{fig:tradeoff} places the evidence levels side by side without filling
missing measurements. The declared \CorrectedNumLayers{}-core BF16 state is
\CorrectedBfSixteenAllLayerStateMiB{} MiB and requires an ideal minimum of
\CorrectedBfSixteenIdealUram{} URAM288 primitives, already above the U55C total
of \CorrectedDeviceAvailableUram{}. Uniform MXFP4 is
\CorrectedUniformMxfpFourAllLayerStateMiB{} MiB and has an ideal lower bound of
\CorrectedUniformMxfpFourIdealUram{} URAM288 plus
\CorrectedUniformMxfpFourIdealBram{} BRAM18K, so it passes only the raw-bit
necessary condition. These totals cover one sequence of recurrence-core state;
they exclude weights, convolution state, full-attention key/value caches,
activations, ports, banking, replication, scratch space, placement, and
routing. Consequently, they do not establish which complete model fits on the
card. The matched all-layer BF16 synthesis requires
\CorrectedBfSixteenRambRequired{} RAMB36/FIFO primitives versus
\CorrectedBfSixteenRambAvailable{} available, so physical fit is
\CorrectedBfSixteenPhysicalFit{} and no BF16 routed timing or power result
exists. A reduced-layer design is not substituted because it would change the
controlled state layout. The HLS panel contains matched BF16, native encoded MXFP4, and native
MXFP8 kernels; matched INT4 remains unrun. No energy axis is shown because
energy was not measured.

\begin{figure*}[t]
\centering
\includegraphics[width=\textwidth]{paper/figures/corrected/tradeoff_evidence.pdf}
\caption{Evidence-level storage, quality, and schedule view, not a measured
Pareto frontier. (a) Logical recurrent-state payload for
\CorrectedNumLayers{} cores; raw PASS/FAIL is only a necessary primitive-bit
capacity condition. (b) Layer-level synthetic output cosine at token
\CorrectedLongTraceTokens{}; floating Q/DQ and native encoded paths remain
distinct. (c) Matched fixed-geometry HLS maximum STEP estimates; the unavailable
INT4 hardware point is marked not run rather than imputed.}
\label{fig:tradeoff}
\end{figure*}

\subsection{Recurrence-Aware Correction}
The correction passes all \CorrectedHeldoutSeedBlocks{} test seed blocks under
both initial-state conditions. Table~\ref{tab:corrected-quality} reports the
worst of the \CorrectedHeldoutPairedConditions{} paired test conditions beside
the two fully recomputed \CorrectedLongTraceTokens{}-token development
conditions. Passing the three-hard-counter gate does not imply event-free
arithmetic: the longer runs contain large alignment-underflow and deliberate
residual-clip counts. The seed block, not each initial-state condition, is the
stochastic unit. These are finite-trace engineering checks, not a
lossless-arithmetic or robustness claim.

\begin{table}[t]
\centering
\caption{Corrected-candidate test and development results. Test values are the
worst across \CorrectedHeldoutSeedBlocks{} seed blocks, each paired with random
and zero initial state. Development values are the worst or total across
\CorrectedExtendedRunCount{} fully recomputed conditions. Dashes denote
diagnostics not aggregated for the test set.}
\label{tab:corrected-quality}
\scriptsize
\resizebox{\columnwidth}{!}{\begin{tabular}{lrr}
\toprule
Metric & Test (\CorrectedHeldoutTraceTokens{}) & Development (\CorrectedLongTraceTokens{}) \\
\midrule
Min. checkpoint cosine & 0.994892 & 0.995065 \\
Min. all-token cosine & 0.994430 & 0.994389 \\
Max. final state rel. $L_2$ & 0.097368 & 0.097472 \\
Max. all-token state rel. $L_2$ & 0.098653 & 0.099057 \\
Max. state abs. error & 0.260358 & 0.267635 \\
Hard-event total (three counters) & 0 & 0 \\
\midrule
Alignment underflows & -- & 1{,}049{,}355{,}823 \\
E2M0 residual clips & -- & 23{,}296{,}266 \\
Log folds & -- & 2{,}340 \\
\bottomrule
\end{tabular}
}
\end{table}

The numerical recovery does not produce an FPGA Pareto improvement.
Table~\ref{tab:mitigation-cost} compares logical storage and HLS estimates with
both baselines. The correction is larger and slower even before physical
implementation; its periodic fold also creates a large command burst. Three
explicit II=\CorrectedTargetIi{} constraints fail.

\begin{table}[t]
\centering
\caption{Logical storage and per-layer HLS command estimates for the
recurrence-aware mitigation. ``Amortized'' charges one fold every
\CorrectedFoldPeriodTokens{} STEPs; it is not service p99.}
\label{tab:mitigation-cost}
\scriptsize
\resizebox{\columnwidth}{!}{\begin{tabular}{lrrr}
\toprule
Metric & BF16 & Encoded MXFP4 & Corrected candidate \\
\midrule
Logical bytes/layer & 1{,}048{,}576 & 278{,}528 & 537{,}744 \\
LUT & 84{,}919 & 138{,}969 & 435{,}089 \\
Non-fold STEP max & 4{,}225{,}728 & 6{,}070{,}592 & 17{,}954{,}144 \\
Amortized per-layer HLS cycles/STEP & 4{,}225{,}728 & 6{,}070{,}592 & 39{,}947{,}365 \\
Fold-command burst max & -- & -- & 153{,}952{,}544 \\
\bottomrule
\end{tabular}
}
\end{table}

\subsection{Post-Route and Generated-RTL Evidence}
\label{sec:postroute}
The corrected candidate physically fits the U55C in the out-of-context flow.
It uses \CorrectedPostrouteLut{} CLB LUTs, \CorrectedPostrouteFf{} flip-flops,
\CorrectedPostrouteBram{} BRAM tiles, \CorrectedPostrouteUram{} URAMs, and
\CorrectedPostrouteDsp{} DSPs. URAM is the limiting reported resource at
\CorrectedPostrouteUramPercent{} percent. The implementation has
\CorrectedPostrouteDrcWarnings{} design-rule warnings but no critical warnings
or errors.

Physical fit does not imply target timing. At
\CorrectedPostrouteTargetMhz{} MHz the design has
\CorrectedPostrouteTwoFiftyWnsNs{} ns WNS, and the
\CorrectedPostrouteSecondaryMhz{}-MHz check also fails. The first passing point
in the tested fixed-route sweep is \CorrectedPostrouteFirstClosingMhz{} MHz.
The \CorrectedPostrouteCriticalPathNs{} ns target-clock critical path is
\CorrectedPostrouteCriticalRoutePercent{} percent routing with
\CorrectedPostrouteCriticalLogicLevels{} logic levels and terminates at a
resident-state URAM reset/write-enable pin. It is not an E2M1 multiplier path.

Official Vitis HLS/XSIM co-simulation passes
\CorrectedEZeroControlCosimCommands{} of
\CorrectedEZeroControlCosimCommands{} generated-RTL control commands at
\CorrectedEZeroControlCosimLatency{} cycles per command. Both commands return
before a recurrent update, so the candidate's
\CorrectedStatefulCsimTokens{}-token exact recurrent result remains C
simulation, not RTL parity. The direct Verilator snapshot LOAD is
\CorrectedEZeroRtlDirectLoadStatus{}, but the harness completes
\CorrectedEZeroRtlCompletedSteps{} recurrent
STEPs; the required 64-token candidate result is therefore
\CorrectedEZeroRtlSixtyFourStatus{}. Table~\ref{tab:postroute} preserves this
boundary.

\begin{table}[t]
\centering
\caption{Corrected-candidate out-of-context physical and generated-RTL
evidence. Power is vectorless; the closing frequency is the first passing
tested fixed-route point.}
\label{tab:postroute}
\scriptsize
\resizebox{\columnwidth}{!}{\begin{tabular}{lrl}
\toprule
Evidence item & Result & Status \\
\midrule
\multicolumn{3}{l}{\textit{Corrected MXFP4 candidate}} \\
Out-of-context physical fit & yes & PASS \\
250-MHz setup WNS & -1.540 ns & FAIL \\
First tested closing point & 180.18 MHz & PASS \\
CLB LUT / FF & 208{,}523 / 111{,}027 & FIT \\
BRAM tile / URAM / DSP & 353 / 624 / 44 & FIT \\
DRC warnings / critical / errors & 38 / 0 / 0 & PASS\_WITH\_WARNINGS \\
Vectorless power at 5.6 ns & 6.991 W & Medium estimate \\
Generated-RTL control smoke & 2 / 2 commands & PASS \\
Candidate 64-token HLS C parity & 64 tokens & PASS \\
Candidate generated-RTL direct LOAD & 1 command & PASS \\
Candidate recurrent 64-token RTL parity & 0 / 64 STEP commands & NOT\_ESTABLISHED \\
\midrule
\multicolumn{3}{l}{\textit{Native MXFP8 baseline}} \\
Out-of-context physical fit & yes & PASS \\
250-MHz setup WNS & -1.355 ns & FAIL \\
First tested closing point & 166.67 MHz & PASS \\
CLB LUT / FF & 58{,}582 / 39{,}807 & FIT \\
BRAM tile / URAM / DSP & 164 / 576 / 2 & FIT \\
Vectorless power at 6.00 ns & 4.952 W & Medium estimate \\
\bottomrule
\end{tabular}
}
\end{table}

At a timing-closed \CorrectedPostroutePowerPeriodNs{} ns period, Vivado reports
\CorrectedPostroutePowerTotalW{} W total on-chip power:
\CorrectedPostroutePowerDynamicW{} W dynamic and
\CorrectedPostroutePowerStaticW{} W static, with
\CorrectedPostroutePowerConfidence{} confidence. No measured activity, board
telemetry, matched routed BF16 design, or recurrent RTL service latency is
available, so we do not derive energy per token or an energy-efficiency ratio.
The matched native-MXFP8 image also fits, using
\CorrectedMxfpEightPostrouteLut{} LUTs and
\CorrectedMxfpEightPostrouteUram{} URAMs. It fails the 250-MHz check and first
closes at \CorrectedMxfpEightPostrouteFirstClosingMhz{} MHz in the tested
fixed-route sweep. Its \CorrectedMxfpEightPostroutePowerTotalW{} W vectorless
estimate is not measured energy.

\section{Discussion}
The controlled answer is no for the current implementations. Floating Q/DQ
already predicts substantial four-bit drift, while the native encoded path is
worse and records \CorrectedNativeEncodedAlignmentUnderflowsEightK{} alignment
underflows through token \CorrectedLongTraceTokens. The encoded HLS point also
costs more than BF16. Adding a dense residual and recent-write log restores the
finite-trace engineering criteria, but its worst longer-run state error is close
to the boundary and its computation costs more than BF16.

This result separates representation density from implementation cost. Shared
scales reduce bits, but exponent comparison, alignment, requantization, and
folding add control and arithmetic. The E2M1 multiplier is small and DSP-free;
the surrounding recurrence is not. Table~\ref{tab:resource-ablation} localizes
the dominant generated hierarchy. The correction's fold engine and update
quantizer account for \CorrectedCorrectionMachineryLutPercent{} percent of its
top-level LUT estimate. A useful replacement must reduce the full state-update
path, not only the encoded element width.

\begin{table}[t]
\centering
\caption{Generated HLS hierarchy ablation. Shares are relative to each
variant's own top-level LUT estimate and are not routed utilization.}
\label{tab:resource-ablation}
\scriptsize
\resizebox{\columnwidth}{!}{\begin{tabular}{lrr}
\toprule
Generated hierarchy block & LUT & Share of top (\%) \\
\midrule
BF16 state-update block & 2{,}191 & 2.6 \\
Encoded-MXFP4 state-update block & 43{,}843 & 31.5 \\
Correction fold engine & 185{,}073 & 42.5 \\
Correction update quantizer & 157{,}419 & 36.2 \\
\midrule
Correction subtotal & 342{,}492 & 78.7 \\
\bottomrule
\end{tabular}
}
\end{table}

The word \emph{drift} is deliberate. For one head, let $E_t$ be candidate minus
reference state. For the same external token inputs, define $R_t$ as the exact
residual between the implemented candidate transition and the ideal reference
transition applied to the candidate state. This residual includes input and
state quantization, alignment, accumulation, scale, clipping, and fold error.
The recurrence then gives
\begin{equation}
E_t=\alpha_t(I-\beta_t k_tk_t^{\mathsf T})E_{t-1}+R_t.
\end{equation}
For normalized $k_t$ and $\beta_t\in[0,1]$,
\begin{equation}
\lVert E_t\rVert_F\leq
\alpha_t\lVert E_{t-1}\rVert_F+\lVert R_t\rVert_F.
\end{equation}
Let $P_t=k_tk_t^{\mathsf T}$. The parallel and orthogonal components obey
the exact identity
\begin{equation}
\begin{split}
\lVert(I-\beta_tP_t)E\rVert_F^2={}&
(1-\beta_t)^2\lVert k_t^{\mathsf T}E\rVert_2^2\\
&+\lVert(I-P_t)E\rVert_F^2.
\end{split}
\end{equation}
Thus the delta gate attenuates only error parallel to the current key; the
worst-case operator norm is one because $K>1$. Unrolling the one-step bound
gives
\begin{equation}
\lVert E_T\rVert_F\leq
\Bigl(\prod_{i=1}^{T}\alpha_i\Bigr)\lVert E_0\rVert_F+
\sum_{j=1}^{T}\Bigl(\prod_{i=j+1}^{T}\alpha_i\Bigr)\lVert R_j\rVert_F.
\end{equation}
If $\alpha_i\leq\bar\alpha<1$ and $\lVert R_i\rVert_F\leq\bar R$, the
right-hand side is at most
$\bar\alpha^T\lVert E_0\rVert_F+
(1-\bar\alpha^T)\bar R/(1-\bar\alpha)$.
Our protocol permits decay arbitrarily close to one, so it supplies no uniform
contraction margin. Changing keys may attenuate different directions, but a
persistent-excitation bound is not established. These relations explain why
repeated local errors matter, but they are not a formal stability certificate. We do
not have a conservative bound on $R_t$ spanning
input and base quantization, logged coefficients, folding, alignment,
accumulation, scale changes, and clipping. ``Stable'' in the engineering gate
therefore means only that the measured finite traces stay within its thresholds.

The correction also changes the experimental role of MXFP8. Uniform MXFP8
state is a stronger numerical baseline than MXFP4, while the corrected
candidate's logical payload is only \CorrectedPercentBelowMxfpEight{} percent
smaller. Both the corrected candidate and native MXFP8 physically fit out of
context, but neither closes 250 MHz. Without measured activity, a routed BF16
comparator, or board energy, their power estimates do not establish an energy
advantage for the more complex representation.

\section{Limitations and Threats to Validity}
The long-horizon numerical evidence is layer-level and synthetic. The short
Qwen diagnostic reconstructs query, key, value, decay, and delta-gate inputs
at one pinned layer boundary, but contains only
\CorrectedQwenRecurrentMinTokens{}--\CorrectedQwenRecurrentMaxTokens{} valid
tokens per prompt. It is not closed-loop execution and supports neither
perplexity nor downstream-task claims. The kernel excludes projections,
normalization, full-attention layers, and their key-value cache.

The nominal study uses one development seed. The correction uses a separate
high-retention trace family and only \CorrectedHeldoutSeedBlocks{} test seed
blocks crossed with two initial states, so it is not a paired numerical
ablation. Its registration was not externally timestamped with the full
dependency closure; replay checks integrity, not independence. Metrics flatten
heads and have no trace-population confidence interval or FP64 long-horizon
curve. Thresholds are not calibrated to model quality. The longer-run error is
only
\CorrectedExtendedStateRelLTwoGateMargin{} below its boundary, and changing that
boundary alone to \CorrectedStateRelLTwoSensitivityProbe{} rejects both test and
development conditions. Section~VII supplies no formal stability certificate.

BF16 and uniform-MXFP4 results remain HLS estimates. The correction and native
MXFP8 baseline are routed out of context; both fit but fail 250 MHz. The
fixed-route sweeps are not global frequency searches, and matched all-layer
BF16 fails capacity before routing. Xilinx Runtime (XRT), xclbin, board checks, and synchronized
same-boundary GPU measurements are absent; vectorless power is not board energy.

Packing and token quantization are excluded, and no multi-layer service latency
is measured. Verification is asymmetric: uniform MXFP4 has a 64-token RTL
trace; the correction has 64-token C parity, a two-command RTL control smoke,
and one exact direct-RTL LOAD. Two isolated XSIM runs exhaust host memory before
transaction one, and no candidate recurrent STEP completes. Results cover one
U55C and one parallelism point; cross-device and longer real-trace studies
remain open.

\section{Conclusion}
Native encoded MXFP4 does not replace BF16 in this persistent-state GDN core:
it has more drift and higher fixed-geometry HLS cost. A residual/write-log
mitigation meets bounded synthetic criteria but exceeds BF16 cost, misses II
and routed clock targets, and lacks recurrent RTL parity. The corrected and
MXFP8 images fit out of context, but longer real traces, closed-loop quality,
and measured board energy remain necessary.

\begin{thebibliography}{00}
\bibitem{gateddeltanet}
S. Yang, J. Kautz, and A. Hatamizadeh, ``Gated Delta Networks: Improving
Mamba2 with Delta Rule,'' in \emph{Proc. Int. Conf. Learning Representations},
2025.

\bibitem{gupta2026}
N. Gupta, P. Wang, R. Kannan, and V. K. Prasanna, ``A Persistent-State
Dataflow Accelerator for Memory-Bound Linear Attention Decode on FPGA,''
arXiv preprint arXiv:2603.05931v1 [cs.AR], Mar. 2026, doi:
10.48550/arXiv.2603.05931. [Online]. Available:
\url{https://arxiv.org/abs/2603.05931}

\bibitem{ocpmx}
Open Compute Project, ``OCP Microscaling Formats (MX) Specification,'' ver.
1.0, Sep. 2023. [Online]. Available:
\url{https://www.opencompute.org/documents/ocp-microscaling-formats-mx-v1-0-spec-final-pdf}

\bibitem{qwen3nextmodel}
Qwen Team, ``Qwen3-Next-80B-A3B-Instruct,'' Hugging Face model repository,
rev. 9c7f2fbe84465e40164a94cc16cd30b6999b0cc7, accessed Aug. 1, 2026.
[Online]. Available:
\url{https://huggingface.co/Qwen/Qwen3-Next-80B-A3B-Instruct/tree/9c7f2fbe84465e40164a94cc16cd30b6999b0cc7}

\bibitem{transformersqwen}
Hugging Face, ``Qwen3-Next implementation,'' \emph{Transformers}, ver. 4.57.0,
commit 8ac2b916b042b1f78b75c9eb941c0f5d2cdd8e10, 2025. [Online]. Available:
\url{https://github.com/huggingface/transformers/blob/8ac2b916b042b1f78b75c9eb941c0f5d2cdd8e10/src/transformers/models/qwen3_next/modeling_qwen3_next.py}

\bibitem{fla021}
FLA Contributors, ``Gated delta rule implementation,''
\emph{flash-linear-attention}, ver. 0.2.1, commit
a670dff4c2537fc1a82486584dd9569e18fba833, 2025. [Online]. Available:
\url{https://github.com/fla-org/flash-linear-attention/tree/a670dff4c2537fc1a82486584dd9569e18fba833/fla/ops/gated_delta_rule}

\bibitem{deltanet}
S. Yang, B. Wang, Y. Zhang, Y. Shen, and Y. Kim, ``Parallelizing Linear
Transformers with the Delta Rule over Sequence Length,'' in \emph{Advances in
Neural Information Processing Systems}, vol. 37, 2024, doi:
10.52202/079017-3668.

\bibitem{qllama}
H. Wen, Z. Wang, J. Dong, W. Lou, L. Gong, C. Wang, and X. Zhou,
``QLlama: An FPGA-Based Microscaling Quantization Accelerator for
Energy-Efficient Llama2 Inference,'' \emph{IEEE Embedded Systems Letters},
vol. 17, no. 5, pp. 337--340, 2025, doi: 10.1109/LES.2025.3600563.

\bibitem{jackscales}
M. Mekhemer, A. Elsousy, B. Venkatesh, R. Rowley, V. Betz, N. Kapre, and
A. Boutros, ``Jack of All Scales: A Versatile FPGA Tensor Block for MXFP
Precisions,'' arXiv preprint arXiv:2607.13898v1 [cs.AR], Jul. 2026, doi:
10.48550/arXiv.2607.13898.

\bibitem{mxformer}
G. Karfakis \emph{et al.}, ``MXFormer: A Microscaling Floating-Point
Charge-Trap Transistor Compute-in-Memory Transformer Accelerator,'' arXiv
preprint arXiv:2602.12480v1 [cs.AR], Feb. 2026, doi:
10.48550/arXiv.2602.12480.

\bibitem{micromix}
W. Liu, H. Meng, Y. Luo, P. Zhang, and X. Ma, ``MicroMix: Efficient
Mixed-Precision Quantization with Microscaling Formats for Large Language
Models,'' in \emph{Proc. Int. Conf. Learning Representations}, 2026.

\bibitem{mxattention}
J. Yu \emph{et al.}, ``MXAttention: Data-Free Optimal Scaling and
Pre-Normalization Quantization for MXFP4 Attention,'' arXiv preprint
arXiv:2607.24377v1, Jul. 2026, doi: 10.48550/arXiv.2607.24377.

\bibitem{mxfpbenchmark}
M. Zhang \emph{et al.}, ``Benchmarking Post-Training Quantization of Large
Language Models under Microscaling Floating Point Formats,'' arXiv preprint
arXiv:2601.09555v1, Jan. 2026, doi: 10.48550/arXiv.2601.09555.

\bibitem{quantgdn}
L. Zhang \emph{et al.}, ``When Good Enough Is Optimal: Multiplication-Only
Matrix Inversion Approximation for Quantized Gated DeltaNet,''
in \emph{Proc. 43rd Int. Conf. Machine Learning}, PMLR, vol. 306, 2026.

\bibitem{kivi}
Z. Liu, J. Yuan, H. Jin, S. Zhong, Z. Xu, V. Braverman, B. Chen, and X. Hu,
``KIVI: A Tuning-Free Asymmetric 2bit Quantization for KV Cache,'' in
\emph{Proc. 41st Int. Conf. Machine Learning}, PMLR, vol. 235,
pp. 32332--32344, 2024.

\bibitem{ssdi8}
H. Kim, B. Ko, M. Kang, M. Kim, D. Lee, J. Lee, S. Yoon, and D. Jung,
``SSDi8: Accurate and Efficient 8-bit Quantization for State Space Duality,''
in \emph{Proc. Int. Conf. Learning Representations}, 2026.
\end{thebibliography}

\end{document}
